% !TEX root = ../main.tex
\begin{ouverture}[electrolyse]{10}{Électrolyse et stockage de l'énergie}{Partie 1 · Constitution et transformations de la matière}
{Transformation forcée\enspace·\enspace Électrolyseur\enspace·\enspace Quantités de matière et courant\enspace·\enspace Accumulateurs et dihydrogène}
\objectif{Distinguer transformation spontanée et transformation forcée par un apport d'énergie électrique.}
\objectif{Modéliser et schématiser un électrolyseur ; identifier l'anode et la cathode.}
\objectif{Relier la durée et l'intensité du courant aux quantités de matière formées aux électrodes.}
\objectif{Décrire le fonctionnement d'un accumulateur, en décharge et en charge.}
\objectif{Discuter des enjeux du stockage de l'énergie sous forme chimique.}
\prerequis{Oxydants, réducteurs, demi-équations électroniques, piles (chapitre~8).}
\prerequis{Critère d'évolution spontanée, quotient de réaction (chapitre~7).}
\prerequis{Charge électrique $Q=I\,\Delta t$ ; énergie électrique $E=U\,Q$.}
\end{ouverture}
\essentiel
\essSection{1}{Transformations spontanées et forcées}
\begin{definition}[Transformation forcée]
Un système chimique évolue spontanément dans le sens qui rapproche $Q_\mathrm{r}$ de $K$. Un générateur électrique peut
imposer un courant qui fait évoluer le système \textbf{en sens inverse} : $Q_\mathrm{r}$ s'éloigne alors de $K$. Une telle
transformation, impossible sans apport d'énergie, est dite \textbf{forcée} : c'est une \textbf{électrolyse}. L'énergie
électrique fournie est convertie en énergie chimique, stockée par les produits formés.
\end{definition}
\begin{exemple}
Dans une solution de bromure de cuivre(II), la transformation \ce{Cu^2+(aq) + 2 Br^-(aq) -> Cu(s) + Br2(aq)} a une
constante d'équilibre très petite ($K\approx\num{e-26}$) : elle ne se produit pas spontanément. Dès qu'un générateur
impose un courant suffisant, du cuivre se dépose sur une électrode et du dibrome orangé apparaît autour de l'autre.
\end{exemple}
\essSection{2}{L'électrolyseur}
\begin{definition}[Électrolyseur, anode et cathode]
Un électrolyseur comporte deux électrodes plongées dans un électrolyte et reliées à un générateur. Comme dans une
pile, l'\textbf{anode} est l'électrode où se produit l'oxydation et la \textbf{cathode} celle où se produit la
réduction. Mais les polarités sont inversées :
\begin{itemize}[leftmargin=1.4em,itemsep=0.1em]
\item l'anode est reliée à la \textbf{borne +} du générateur, qui lui arrache des électrons ;
\item la cathode est reliée à la \textbf{borne $-$}, qui lui fournit des électrons.
\end{itemize}
\end{definition}
\schema{10-electrolyseur}
\begin{savoirfaire}[Analyser une électrolyse]
Repérer les bornes du générateur pour identifier l'anode et la cathode. Écrire l'oxydation à l'anode et la réduction à
la cathode, parmi les couples des espèces présentes (ions, solvant, électrodes). En déduire l'équation de
l'électrolyse et les observations attendues.
\end{savoirfaire}
\begin{attention}
Dans une pile, l'anode est le pôle $-$ ; dans un électrolyseur, elle est reliée au pôle $+$ du générateur. Dans les deux
cas, retenir : \textbf{anode = oxydation}, \textbf{cathode = réduction}, et le courant entre dans l'électrolyseur par
l'anode.
\end{attention}
\essSection{3}{Quantités de matière et quantité d'électricité}
\begin{propriete}[Relier courant et quantités de matière]
Pendant une durée $\Delta t$, un courant d'intensité constante $I$ fait circuler la charge $Q=I\,\Delta t$, soit une
quantité d'électrons
\[n(\mathrm{e}^-)=\frac{Q}{F}=\frac{I\,\Delta t}{F},\qquad F=\qty{96500}{C.mol^{-1}}.\]
Les demi-équations relient ensuite $n(\mathrm{e}^-)$ aux quantités de matière transformées à chaque électrode.
\end{propriete}
\begin{exemple}
Pour cuivrer une pièce, on la place à la cathode d'un électrolyseur contenant des ions \ce{Cu^2+}, avec $I=\qty{1,2}{A}$
pendant 30~min : $n(\mathrm{e}^-)=\frac{1{,}2\times1800}{\num{96500}}=\qty{2,24e-2}{mol}$. La réduction
\ce{Cu^2+ + 2 e- -> Cu} consomme deux électrons par atome : $n(\ce{Cu})=\qty{1,12e-2}{mol}$, soit une masse déposée
$m=\num{1.12e-2}\times63{,}5=\qty{0,71}{g}$.
\end{exemple}
\begin{codecote}[0.58]
\begin{pythoncol}
F, M = 96500, 63.5     # C/mol, g/mol
I = 1.2                # A
def masse(t):          # t en secondes
    n_e = I * t / F
    return M * n_e / 2
for t in (600, 1800, 3600):
    print(t, round(masse(t), 2))
\end{pythoncol}
\begin{consolecol}
600 0.24
1800 0.71
3600 1.42
\end{consolecol}
\cote
\fontsize{9.5}{12.5}\selectfont
\capacite{Capacité numérique}\par\smallskip
La fonction calcule la masse de cuivre déposée après une durée $t$ : la quantité d'électrons $I\,t/F$, puis la masse
$M\,n(\mathrm{e}^-)/2$, deux électrons étant échangés par atome. À intensité constante, la masse déposée est
proportionnelle à la durée : on retrouve les 0,71~g de l'exemple.
\end{codecote}
\essSection{4}{Stocker l'énergie sous forme chimique}
\begin{definition}[Accumulateur]
Un \textbf{accumulateur} est une pile rechargeable : en \textbf{décharge}, il fonctionne en pile et convertit son
énergie chimique en énergie électrique ; en \textbf{charge}, un générateur force la transformation inverse et
reconstitue les réactifs. À chaque électrode, la demi-équation s'inverse entre les deux modes : l'électrode qui était la
cathode en décharge devient l'anode en charge.
\end{definition}
\schema{10-accumulateur}
\begin{propriete}[Le stockage chimique de l'énergie]
L'électricité se stocke mal directement : on la convertit en énergie chimique, dans des accumulateurs (téléphones,
véhicules électriques) ou dans du dihydrogène produit par électrolyse de l'eau, converti ensuite par une pile à
combustible. Chaque conversion s'accompagne de pertes : le rendement global d'une chaîne de stockage est le produit des
rendements de chaque étape.
\end{propriete}
\begin{experience}[Électrolyser l'eau au laboratoire]
Deux électrodes inertes plongent dans de l'eau rendue conductrice par un peu d'hydroxyde de sodium ; deux éprouvettes
retournées, remplies de solution, coiffent les électrodes. Dès que le générateur débite, des bulles apparaissent : le
volume de gaz recueilli à la cathode est double de celui recueilli à l'anode. Le premier gaz aboie à la flamme
(dihydrogène), le second rallume une bûchette incandescente (dioxygène) : on a réalisé
\ce{2 H2O(l) -> 2 H2(g) + O2(g)}, transformation forcée.
\end{experience}
\begin{rediger}
Toujours justifier l'anode et la cathode par la borne du générateur à laquelle l'électrode est reliée, puis par la
nature de la réaction. Dans un calcul de quantité d'électricité, convertir la durée en secondes et conclure avec les
unités : « $n(\mathrm{e}^-)=I\,\Delta t/F$, avec $I$ en ampères et $\Delta t$ en secondes ».
\end{rediger}
\exercices

\rubrique{Découverte}
\begin{blocLivre}{ex-10-1}
\begin{exercice}[1]{Pile ou électrolyseur ?}
Deux béchers contiennent le même système chimique : une lame de cuivre dans une solution de sulfate de cuivre(II) et
une lame de zinc dans une solution de sulfate de zinc, reliées par un pont salin. Dans le montage A, les lames sont
reliées à un conducteur ohmique ; dans le montage B, à un générateur dont la borne $+$ est connectée à la lame de
cuivre et qui impose un courant inverse de celui du montage A.
\begin{questions}
\item Dans quel montage le système évolue-t-il spontanément ? Écrire l'équation de la transformation correspondante.
\item Comment le quotient de réaction évolue-t-il par rapport à $K$ dans chaque montage ?
\item Écrire l'équation de la transformation du montage B. Pourquoi la qualifie-t-on de forcée ?
\item Identifier l'anode dans chaque montage.
\end{questions}
\end{exercice}
\end{blocLivre}
\begin{blocLivre}{ex-10-2}
\begin{exercice}[1]{Les électrodes d'un électrolyseur}
On électrolyse une solution de bromure de cuivre(II), $(\ce{Cu^2+}+2\,\ce{Br-})$, entre deux électrodes de graphite.
Au bout de quelques minutes, un dépôt rougeâtre couvre l'électrode 1 et une coloration orangée se développe autour de
l'électrode 2.
\schema{10-ex02}
\begin{questions}
\item Interpréter les deux observations : quelles espèces se forment ?
\item Écrire la demi-équation de la réaction à chaque électrode, à partir des couples \ce{Cu^2+}/\ce{Cu} et
\ce{Br2}/\ce{Br-}.
\item Identifier l'anode et la cathode, puis les bornes du générateur.
\item Écrire l'équation de l'électrolyse. Que se passerait-il si l'on débranchait le générateur pour le remplacer par
un conducteur ohmique ?
\end{questions}
\end{exercice}
\end{blocLivre}
\begin{blocLivre}{ex-10-3}
\begin{exercice}[1]{Quantité d'électricité et quantité de matière}
Un électrolyseur est parcouru par un courant d'intensité constante $I=\qty{0,50}{A}$ pendant 20~min.
\begin{donnees}
$F=\qty{96500}{C.mol^{-1}}$ ; $N_\mathrm{A}=\qty{6,02e23}{mol^{-1}}$.
\end{donnees}
\begin{questions}
\item Calculer la charge électrique qui a traversé l'électrolyseur.
\item En déduire la quantité d'électrons échangés, puis leur nombre.
\item À la cathode se produit la réduction \ce{Ag+ + e- -> Ag}. Quelle quantité d'argent s'est déposée ?
\item Quelle durée faudrait-il, à la même intensité, pour déposer le double ?
\end{questions}
\end{exercice}
\end{blocLivre}
\begin{blocLivre}{ex-10-4}
\begin{exercice}[1]{Cuivrer une médaille}
Pour cuivrer une médaille de surface totale $S=\qty{60}{cm^2}$, on la place à la cathode d'un bain contenant des ions
\ce{Cu^2+}, et l'on fait circuler un courant d'intensité $I=\qty{0,60}{A}$ pendant 40~min.
\begin{donnees}
$M(\ce{Cu})=\qty{63,5}{g.mol^{-1}}$ ; masse volumique du cuivre : $\rho=\qty{8,96}{g.cm^{-3}}$ ;
$F=\qty{96500}{C.mol^{-1}}$.
\end{donnees}
\begin{questions}
\item Écrire la demi-équation de la réaction sur la médaille. Justifier qu'elle se produit à la cathode.
\item Calculer la quantité d'électrons échangés, puis la masse de cuivre déposée.
\item En déduire l'épaisseur moyenne du dépôt, en micromètres.
\item Comment doubler cette épaisseur sans changer l'intensité ?
\end{questions}
\end{exercice}
\end{blocLivre}
\begin{blocLivre}{ex-10-5}
\begin{exercice}[1]{Recharger un accumulateur}
Une pile Daniell usée est branchée sur un générateur : la borne $+$ du générateur est reliée à l'électrode de cuivre,
et le courant imposé est inverse de celui que la pile débitait.
\begin{questions}
\item Écrire la réaction qui se produisait à chaque électrode pendant la décharge.
\item Écrire les réactions aux électrodes pendant la charge. Que devient le dépôt de cuivre ?
\item Pendant la charge, quelle électrode est l'anode ? Comparer à la décharge.
\item Pourquoi la charge est-elle une transformation forcée ? D'où provient l'énergie stockée ?
\end{questions}
\end{exercice}
\end{blocLivre}
\rubrique{Entraînement}
\begin{blocLivre}{ex-10-6}
\begin{exercice}[2]{L'électrolyse de l'eau}
Deux électrodes inertes plongent dans de l'eau rendue conductrice ; deux éprouvettes graduées, retournées et remplies
de solution, recueillent les gaz. Le générateur débite $I=\qty{0,25}{A}$ pendant 15,0~min.
\schema{10-ex06}
\begin{donnees}
Couples : \ce{O2}/\ce{H2O} et \ce{H2O}/\ce{H2} (\ce{2 H2O + 2 e- <=> H2 + 2 HO-}) ; volume molaire :
$V_\mathrm{m}=\qty{24}{L.mol^{-1}}$ ; $F=\qty{96500}{C.mol^{-1}}$.
\end{donnees}
\begin{questions}
\item Écrire la demi-équation de la réaction à chaque électrode, puis l'équation de l'électrolyse.
\item Calculer la quantité d'électrons échangés pendant l'expérience.
\item En déduire les volumes de dihydrogène et de dioxygène recueillis. Dans quel rapport sont-ils ?
\item Comment identifier chacun des deux gaz ?
\end{questions}
\end{exercice}
\end{blocLivre}
\begin{blocLivre}{ex-10-7}
\begin{exercice}[2]{Argenter des couverts}
Un orfèvre argente des couverts : les pièces, placées à la cathode d'un bain d'ions \ce{Ag+}, doivent recevoir 2,0~g
d'argent au total. L'anode est une lame d'argent pur. L'intensité vaut $I=\qty{0,50}{A}$.
\begin{donnees}
$M(\ce{Ag})=\qty{107,9}{g.mol^{-1}}$ ; $F=\qty{96500}{C.mol^{-1}}$.
\end{donnees}
\begin{questions}
\item Écrire la demi-équation de la réaction sur les couverts.
\item Calculer la quantité d'électrons nécessaire, puis la durée de l'opération.
\item À l'anode, l'argent de la lame est oxydé. Écrire la demi-équation correspondante.
\item Montrer que la concentration en ions \ce{Ag+} du bain reste constante. Quel est l'intérêt de cette « anode
soluble » ?
\end{questions}
\end{exercice}
\end{blocLivre}
\begin{blocLivre}{ex-10-8}
\begin{exercice}[2]{L'affinage du cuivre}
Pour purifier le cuivre issu des minerais (pureté 99~\%), on en fait l'anode d'un électrolyseur ; la cathode est une
fine feuille de cuivre pur. L'anode se dissout, le cuivre se dépose pur sur la cathode et les impuretés tombent au
fond de la cuve. Une cuve industrielle fonctionne sous 0,30~V avec $I=\qty{250}{A}$.
\begin{donnees}
$M(\ce{Cu})=\qty{63,5}{g.mol^{-1}}$ ; $F=\qty{96500}{C.mol^{-1}}$ ; $\qty{1}{kW.h}=\qty{3,6e6}{J}$.
\end{donnees}
\begin{questions}
\item Écrire les demi-équations aux deux électrodes. Pourquoi la concentration en ions \ce{Cu^2+} ne varie-t-elle
pas ?
\item Calculer la masse de cuivre pur déposée en 24~h de fonctionnement.
\item Calculer l'énergie électrique consommée en 24~h, en kilowattheures, puis par kilogramme de cuivre affiné.
\item Pourquoi ce procédé serait-il inutilisable pour purifier un métal plus réducteur que l'eau, comme le sodium ?
\end{questions}
\end{exercice}
\end{blocLivre}
\begin{blocLivre}{ex-10-9}
\begin{exercice}[2]{L'accumulateur au plomb}
La batterie de démarrage d'une voiture est un accumulateur au plomb. En décharge, les réactions aux électrodes sont :
\[\ce{Pb + SO4^2- -> PbSO4 + 2 e-}\qquad\text{et}\qquad\ce{PbO2 + 4 H+ + SO4^2- + 2 e- -> PbSO4 + 2 H2O}.\]
La batterie a une tension de 12~V et une capacité de \qty{60}{A.h}.
\begin{donnees}
$M(\ce{Pb})=\qty{207,2}{g.mol^{-1}}$ ; $F=\qty{96500}{C.mol^{-1}}$.
\end{donnees}
\begin{questions}
\item Identifier l'anode et la cathode pendant la décharge, puis écrire l'équation de fonctionnement.
\item Convertir la capacité en coulombs et calculer la quantité d'électrons correspondante.
\item Quelle masse de plomb l'anode perd-elle au cours d'une décharge complète ?
\item Le démarreur consomme 150~A pendant 2,0~s. Quelle fraction de la capacité un démarrage utilise-t-il ?
\item Écrire les réactions aux électrodes pendant la charge par l'alternateur.
\end{questions}
\end{exercice}
\end{blocLivre}
\begin{blocLivre}{ex-10-10}
\begin{exercice}[2]{Produire du dihydrogène}
Un électrolyseur de laboratoire industriel, parcouru par un courant de 1500~A, produit du dihydrogène à la cathode :
\ce{2 H2O + 2 e- -> H2 + 2 HO-}. Le programme suivant calcule la masse produite, heure par heure, jusqu'à atteindre
1,0~kg.
\begin{python}
F, I = 96500, 1500           # C/mol, A
M_H2 = 2.0                   # g/mol
t, m = 0, 0
while m < 1000:              # masse en g
    t = t + 3600             # pas d'une heure
    m = M_H2 * I * t / (2 * F)
print(t, round(t / 3600), round(m))
\end{python}
\begin{console}
64800 18 1007
\end{console}
\begin{questions}
\item \capacite{Capacité numérique}Justifier l'expression de \code{m} utilisée par le programme.
\item Interpréter les trois nombres affichés.
\item Retrouver par un calcul direct la durée nécessaire pour produire exactement 1,0~kg de dihydrogène.
\item Quel volume ce kilogramme occuperait-il à 25~°C sous 1,0~bar ($V_\mathrm{m}=\qty{24,8}{L.mol^{-1}}$) ? Comment le
stocke-t-on en pratique ?
\end{questions}
\end{exercice}
\end{blocLivre}
\begin{blocLivre}{ex-10-11}
\begin{exercice}[2]{Protéger l'acier par électrozingage}
Pour protéger une plaque d'acier de la corrosion, on la recouvre de zinc par électrolyse : la plaque, de surface
$S=\qty{0,96}{m^2}$ pour la face traitée, constitue la cathode d'un bain d'ions \ce{Zn^2+}. On vise un dépôt
d'épaisseur $e=\qty{10}{\micro m}$, avec $I=\qty{200}{A}$.
\begin{donnees}
$M(\ce{Zn})=\qty{65,4}{g.mol^{-1}}$ ; $\rho(\ce{Zn})=\qty{7,14}{g.cm^{-3}}$ ; $F=\qty{96500}{C.mol^{-1}}$.
\end{donnees}
\begin{questions}
\item Écrire la demi-équation de la réaction sur la plaque.
\item Calculer le volume, puis la masse de zinc à déposer.
\item En déduire la quantité d'électrons nécessaire et la durée de l'électrolyse.
\item Le zinc est plus réducteur que le fer. Expliquer pourquoi le revêtement protège l'acier même lorsqu'il est
localement rayé.
\end{questions}
\end{exercice}
\end{blocLivre}
\rubrique{Objectif bac}
\begin{blocLivre}{ex-10-12}
\begin{exercice}[3]{Recharger une voiture électrique}
\begin{docu}{1}{La batterie d'une citadine}
La batterie d'une citadine électrique stocke une énergie de 60~kWh. Elle regroupe des cellules lithium-ion de tension
3,7~V et de capacité \qty{5,0}{A.h}. Une prise domestique renforcée délivre une puissance de 11~kW.
\end{docu}
\begin{docu}{2}{La charge d'une cellule lithium-ion}
Pendant la charge, les ions lithium \ce{Li+} quittent l'électrode d'oxyde métallique, traversent l'électrolyte et
s'insèrent dans l'électrode de graphite, où se produit la réaction \ce{Li+ + e- + 6 C -> LiC6}. À la décharge, la
réaction s'inverse.
\end{docu}
\begin{donnees}
$M(\ce{Li})=\qty{6,94}{g.mol^{-1}}$ ; $F=\qty{96500}{C.mol^{-1}}$ ; énergie d'une cellule : $E=U\,Q$.
\end{donnees}
\begin{questions}
\item Calculer l'énergie stockée par une cellule, en wattheures, puis le nombre de cellules de la batterie.
\item Estimer la durée d'une charge complète sur la prise renforcée.
\item Pendant la charge, l'électrode de graphite est-elle l'anode ou la cathode ? À quelle borne du chargeur est-elle
reliée ?
\item \initiative Calculer la masse de lithium qui s'insère dans le graphite d'une cellule au cours d'une charge
complète.
\item La réaction d'insertion s'inverse à la décharge. En quoi cette réversibilité est-elle la propriété essentielle
d'un accumulateur ?
\end{questions}
\end{exercice}
\end{blocLivre}
\begin{blocLivre}{ex-10-13}
\begin{exercice}[3]{Le dihydrogène vert}
\begin{docu}{1}{Un électrolyseur industriel}
Un électrolyseur alcalin industriel regroupe 160 cellules en série, chacune sous une tension de 2,0~V, parcourues par
un courant de 25~kA : sa puissance électrique atteint 8,0~MW. Chaque cellule produit du dihydrogène à sa cathode et du
dioxygène à son anode. Alimenté par de l'électricité renouvelable, il produit un dihydrogène dit « vert ».
\end{docu}
\begin{donnees}
Cathode : \ce{2 H2O + 2 e- -> H2 + 2 HO-} ; anode : \ce{4 HO- -> O2 + 2 H2O + 4 e-} ;
$M(\ce{H2})=\qty{2,0}{g.mol^{-1}}$ ; $F=\qty{96500}{C.mol^{-1}}$ ; énergie libérée par la combustion du dihydrogène :
\qty{33,3}{kW.h} par kilogramme ; $\qty{1}{kW.h}=\qty{3,6e6}{J}$.
\end{donnees}
\begin{questions}
\item Vérifier l'équation de l'électrolyse, \ce{2 H2O -> 2 H2 + O2}, et la valeur de la puissance électrique.
\item Calculer la quantité de dihydrogène produite chaque seconde par une cellule, puis par l'électrolyseur entier.
\item En déduire la masse de dihydrogène produite par jour.
\item \initiative Calculer l'énergie électrique consommée par jour, puis par kilogramme de dihydrogène. Comparer à
l'énergie restituée par sa combustion et conclure sur le rendement du stockage.
\item À quelle condition ce dihydrogène mérite-t-il le qualificatif de « vert » ?
\end{questions}
\end{exercice}
\end{blocLivre}
\begin{blocLivre}{ex-10-14}
\begin{exercice}[3]{De l'aluminium par électrolyse}
\begin{docu}{1}{Le procédé Hall-Héroult}
L'aluminium est produit par électrolyse de l'alumine \ce{Al2O3}, dissoute à 960~°C dans la cryolithe fondue. Les ions
\ce{Al^3+} sont réduits à la cathode ; les anodes, en carbone, sont consommées par la réaction
\ce{C + 2 O^2- -> CO2 + 4 e-}. Une cuve moderne fonctionne sous 4,2~V avec un courant de 350~kA.
\end{docu}
\begin{docu}{2}{Recycler plutôt que produire}
Refondre de l'aluminium recyclé ne demande qu'environ 5~\% de l'énergie nécessaire à sa production par électrolyse.
\end{docu}
\begin{donnees}
$M(\ce{Al})=\qty{27,0}{g.mol^{-1}}$ ; $M(\ce{C})=\qty{12,0}{g.mol^{-1}}$ ; $M(\ce{CO2})=\qty{44,0}{g.mol^{-1}}$ ;
$F=\qty{96500}{C.mol^{-1}}$ ; $\qty{1}{MW.h}=\qty{3,6e9}{J}$.
\end{donnees}
\begin{questions}
\item Écrire la demi-équation de la réaction à la cathode.
\item Calculer la quantité d'électrons échangés par jour, puis la masse d'aluminium produite par une cuve en 24~h.
\item Calculer la masse de carbone consommée et la masse de dioxyde de carbone dégagée chaque jour.
\item \initiative Calculer l'énergie électrique consommée par tonne d'aluminium produite.
\item À l'aide du document~2, expliquer l'intérêt énergétique et climatique du recyclage de l'aluminium.
\end{questions}
\end{exercice}
\end{blocLivre}
\begin{blocLivre}{ex-10-15}
\begin{exercice}[3]{Stocker l'énergie d'une maison autonome}
\begin{docu}{1}{Un surplus d'été à retrouver l'hiver}
Une maison autonome produit, de juin à août, un surplus solaire de 8,0~kWh par jour pendant 90~jours, soit 720~kWh à
stocker pour l'hiver. Deux solutions sont étudiées : des batteries lithium-ion, ou la conversion en dihydrogène.
\end{docu}
\begin{docu}{2}{Deux chaînes de stockage}
Batteries lithium-ion : énergie massique de 150~Wh par kilogramme, rendement charge-décharge de 90~\%, mais
autodécharge et vieillissement sur plusieurs mois. Chaîne dihydrogène : électrolyseur de rendement 70~\%, stockage du
gaz comprimé à 300~bar, pile à combustible de rendement 50~\%.
\end{docu}
\begin{donnees}
Énergie de combustion du dihydrogène : \qty{33,3}{kW.h.kg^{-1}} ; $M(\ce{H2})=\qty{2,0}{g.mol^{-1}}$ ;
$R=\qty{8,314}{J.mol^{-1}.K^{-1}}$ ; $T=\qty{298}{K}$.
\end{donnees}
\begin{questions}
\item Quelle énergie la solution « batteries » restituerait-elle l'hiver ? Calculer la masse de batteries nécessaire
pour stocker les 720~kWh.
\item Calculer le rendement global de la chaîne dihydrogène, puis l'énergie qu'elle restituerait.
\item Calculer la masse de dihydrogène produite à partir des 720~kWh, puis la quantité de matière correspondante.
\item Quel volume ce gaz occupe-t-il à 300~bar et 298~K, en supposant le modèle du gaz parfait applicable ?
\item \initiative Comparer les deux solutions pour un stockage saisonnier. Quel critère, absent des calculs
précédents, favorise le dihydrogène ?
\end{questions}
\end{exercice}
\end{blocLivre}
